\section{Introduction}
\label{sec:intro}

Autoregressive language models built on the Transformer architecture have
become the dominant paradigm in natural language processing, powering
systems such as GPT, LLaMA, and PaLM
\parencite{vaswani2017attention, radford2019language, brown2020language}.
The core innovation of the Transformer is the self-attention mechanism,
which allows each token to attend to every other token in the sequence
while remaining fully parallelisable across the time dimension. In a
\emph{decoder-only} configuration, attention is constrained by a causal
mask so that position $t$ can only attend to positions $\leq t$, making
the model suitable for autoregressive generation where each token is
predicted conditioned on the entire prefix before it.

\noindent The motivation for this work is that implementing a Transformer from
scratch without relying on high-level abstractions such as
\texttt{nn.TransformerDecoder} or \texttt{nn.Multihead} -{Atention} forces
an intimate understanding of the underlying tensor operations, gradient
flow, and numerical stability concerns. These details are precisely the
ones that distinguish a model that merely trains from one that trains
efficiently and generalises well \parencite{karpathy2023nanogpt}. Framed
differently: most practitioners can call a Transformer library function
and get a working model. Far fewer can explain why the causal mask is
necessary for correctness, why the square-root scaling in the attention
formula prevents softmax saturation, or why decaying a LayerNorm gain
hurts. This report is an attempt to close that gap in my own
understanding by building the entire stack from primitives.

\subsection{Problem Statement}
While high-level libraries make it easy to instantiate a Transformer,
they abstract away the causal masking logic, mixed-precision numerics,
and parameter-group weight decay that are essential for competitive
performance. The problem addressed in this report is to implement a
complete, production-style decoder-only Transformer using only PyTorch
tensor primitives, train it on a real text corpus, and verify that its
training and generation behaviour matches the theoretical expectations
set out in the literature.

\subsection{Objectives}
The specific objectives of this assignment are:
\begin{enumerate}
    \item To implement token and sinusoidal positional embeddings from
          scratch.
    \item To implement causal multi-head self-attention manually, including
          the scaled dot-product formulation and the lower-triangular mask.
    \item To build a Pre-LN Transformer block with residual connections and
          a GELU feed-forward network.
    \item To implement mixed-precision training using
          \texttt{torch.autocast} together with a
          \texttt{GradScaler} (active under fp16, disabled under bf16).
    \item To implement a cosine-annealing learning-rate schedule with a
          linear warmup phase.
    \item To apply AdamW with weight decay restricted to 2D weight matrices,
          excluding LayerNorm gains and biases.
    \item To write an autoregressive generation loop supporting greedy,
          temperature, and top-$k$ sampling.
    \item To evaluate training dynamics quantitatively (loss, perplexity)
          and qualitatively (sample generations, attention heatmaps).
\end{enumerate}

\subsection{Summary of Approach}
\label{subsec:approach}

To give the reader a compact view of what was actually built and trained,
Table~\ref{tab:approach-summary} summarises the specific methods,
techniques, and design choices adopted in this work. The remainder of
this report expands on each of them.

\begin{table}[H]
\centering
\caption{Summary of methods and techniques used in this implementation.}
\label{tab:approach-summary}
\small
\begin{tabular}{p{3.2cm} p{10.5cm}}
\toprule
\textbf{Category} & \textbf{Technique / Choice} \\
\midrule
Architecture          & Decoder-only Transformer (GPT-style), $4$ blocks,
                        $d = 128$, $h = 4$ heads, context length $128$,
                        $\approx 0.82$~M parameters. \\
Positional encoding   & Sinusoidal (Vaswani et al.)---chosen over learned
                        embeddings for extrapolation and parameter
                        economy. \\
Attention             & Causal multi-head self-attention implemented
                        manually. Fused QKV projection
                        ($\operatorname{Linear}(d, 3d)$), split into heads,
                        scaled dot-product with a lower-triangular mask
                        registered as a buffer. \\
Block normalisation   & Pre-LayerNorm (LayerNorm before each sub-layer)
                        with residual connections throughout. \\
Feed-forward network  & Two-layer MLP with GELU activation and
                        $4\times$ hidden-dimension expansion. \\
Weight tying          & LM head shares weights with token embedding
                        matrix. \\
Optimiser             & AdamW with $\beta_1 = 0.9$, $\beta_2 = 0.95$.
                        Weight decay $\lambda = 0.1$ applied only to
                        tensors with $\dim \geq 2$ ($18$ tensors); excluded
                        from 1D parameters ($34$ tensors). \\
LR schedule           & Linear warmup ($200$ steps) followed by cosine
                        decay to a floor of $3 \times 10^{-5}$ over
                        $5000$ steps. \\
Mixed precision       & \texttt{torch.autocast} under bf16, with
                        \texttt{GradScaler} disabled (bf16 shares fp32's
                        exponent range). \\
Gradient handling     & Global-norm clipping to $1.0$, applied after
                        unscaling. \\
Generation            & Autoregressive sampling with three modes: greedy
                        decoding, temperature scaling, and top-$k$
                        sampling. \\
Evaluation            & Training/validation loss, perplexity, qualitative
                        samples across five decoding configurations, and
                        attention-weight visualisation. \\
\bottomrule
\end{tabular}
\end{table}

\subsection{Contributions}
This report contributes three things. First, a fully-reproducible,
commented implementation of a GPT-style language model with no black-box
components: every tensor operation, from the embedding lookup to the final
logit computation, is inspectable in the accompanying code. Second, a
systematic evaluation of training dynamics under mixed precision and a
warmup-cosine schedule, including the concrete observation that the
Pre-LN architecture allowed a $200$-step warmup at a peak learning rate
of $3 \times 10^{-4}$ without any early-training instability. Third, a
qualitative study of five sampling strategies at inference time, showing
concretely how the coherence--diversity trade-off described by
\textcite{holtzman2020curious} manifests in a small, character-level
model.
